% P2 (Joanna): report sections 3 (Transistor models), 6 (Error analysis, transistor part)
% and 7 (Limitations and improvements).
%
% Use:   % P2 (Joanna): report sections 3 (Transistor models), 6 (Error analysis, transistor part)
% and 7 (Limitations and improvements).
%
% Use:   % P2 (Joanna): report sections 3 (Transistor models), 6 (Error analysis, transistor part)
% and 7 (Limitations and improvements).
%
% Use:   % P2 (Joanna): report sections 3 (Transistor models), 6 (Error analysis, transistor part)
% and 7 (Limitations and improvements).
%
% Use:   \input{sections/p2_transistor}   from the main report (or pandoc it into report.md).
% Needs: \usepackage{amsmath, array, booktabs, graphicx, xcolor}
%        \newcommand{\todo}[1]{\textcolor{red}{[TODO: #1]}}
%
% Numbers marked "first-round" come from the models committed in 1f1d9f7 / validation
% experiments run on 2026-11-07. Everything that depends on the FINAL retrain (after the
% low-Vds data from P1) is marked \todo and must be pasted from results/tables/*.md.

\section{Transistor models}\label{sec:models}

\subsection{Data split}
Each of the four single-transistor datasets (\texttt{nmos\_on}, \texttt{nmos\_off},
\texttt{pmos\_on}, \texttt{pmos\_off}) is shuffled with a fixed seed (\texttt{seeds.split}~$=42$)
and split 70/15/15 into training, validation and test sets (7\,000/1\,500/1\,500 rows for
10\,000 samples). The row indices of every split are saved in \texttt{artifacts/splits/} so
the split is reproducible and auditable. All model choices---model type, hyperparameters and
preprocessing---are made on the validation set only; the test set is used exactly once, in
\texttt{scripts/05\_finalize\_models.py}, to report the chosen models.
\todo{update counts if P1's low-$V_{DS}$ sampling changes the dataset size}

\subsection{Input features}
The raw inputs are $V_{DD}$, $T$, $W$, $V_D$, $V_S$, $V_G$, $V_B$ and the five varied BSIM4
parameters (\texttt{vth0}, \texttt{toxe}, \texttt{ndep}, \texttt{u0}, \texttt{xj}). Within one
case $V_G$ and $V_B$ are always $0$ or $V_{DD}$, so they carry no information on their own; the
terminal voltage \emph{differences} do. From the raw inputs we build 22 features:
\begin{itemize}
  \item supply and temperature: $V_{DD}$, $T$ (K), $1/T$, thermal voltage $v_T = kT/q$;
  \item geometry: $W$, $\ln W$;
  \item bias: $V_{GS}$, $V_{DS}$, $V_{BS}$, $V_{GD}$, $V_{BD}$, $|V_{DS}|$;
  \item subthreshold physics: $V_{GS}/v_T$, $V_{DS}/v_T$, $(V_{GS}-|V_{th0}|)/v_T$ and the drain
        factor $1-e^{-|V_{DS}|/v_T}$;
  \item process: the five parameters divided by their nominal values from the model card
        (so they lie in $[0.9, 1.1]$), plus $1/t_{ox}$, since gate tunnelling grows steeply as
        the oxide gets thinner.
\end{itemize}
For PMOS every voltage difference is negated (e.g.\ $V_{SG}$ instead of $V_{GS}$) and
$|V_{th0}|$ is used, so PMOS features read like NMOS ones: $V_{GS}>0$ always means ``more on''.
The same function computes the features in training and inside the saved model, and the
feature list is stored in the model file; a model refuses to predict if the feature code
has changed since it was trained.

\subsection{Target transform}
The four terminal currents span about ten decades (from $\sim\!10^{-22}$\,A for OFF-state
gate current to $\sim\!10^{-4}$\,A for ON-state drain current) and $I_D$, $I_S$ change sign
with $V_D-V_S$. Regressing amperes with a squared loss would fit only the largest currents, so
each model predicts
\begin{equation}
  z_t = \operatorname{asinh}\!\left(\frac{I_t}{i_{0,t}}\right),\qquad
  t \in \{I_D, I_G, I_S, I_B\},
\end{equation}
where $i_{0,t}$ is the 1st percentile of $|I_t|$ on the training rows. For $|I|\gg i_0$,
$\operatorname{asinh}(I/i_0)\approx \operatorname{sign}(I)\ln(2|I|/i_0)$, so an error
$\delta z$ is a relative current error of about $\delta z$, which is exactly what we are judged
on. Unlike $\log_{10}|I|$, asinh is smooth through zero and carries the sign, so no separate sign
model is needed. A per-target $i_0$ is necessary because the currents live on very different
scales: one fixed $i_0$ cannot be ``small'' for all of them.

\paragraph{Channel shaping (ON cases).}
In the ON cases the channel current passes through zero at $V_D=V_S$ with the shape
$1-e^{-V_{DS}/v_T}$ in subthreshold and $\propto V_{DS}$ in the linear region. For
\texttt{nmos\_on} and \texttt{pmos\_on}, $I_D$ and $I_S$ are therefore divided by
\begin{equation}
  s(V_{DS}) = \tanh\!\left(\frac{V_D - V_S}{2 v_T}\right)
\end{equation}
before the transform and multiplied back afterwards; the model then learns a smooth,
one-signed quantity. On the validation set this reduced the 95th-percentile error of the
\texttt{nmos\_on} MLP from 164\,\% / 206\,\% to 34\,\% / 48\,\% for $I_D$ / $I_S$. It is
\emph{not} used for the OFF cases: there $I_D$ and $I_S$ near $V_D=V_S$ are dominated by
junction leakage, which does not change sign with $V_{DS}$, and shaping made them worse
(\texttt{nmos\_off} GBDT: 65\,\% / 100\,\% $\rightarrow$ 70\,\% / 111\,\%).

\subsection{Model types and hyperparameters}
Four model types are compared for every case (Table~\ref{tab:hyper}). Each predicts all four
transformed currents at once. Inputs are standardised, and every estimator is seeded
(\texttt{seeds.training}~$=7$). Each model type has a small grid (6 points) that is searched
on the validation set; the grid point with the lowest validation score is kept.

\begin{table}[h]
\centering\small
\begin{tabular}{@{}>{\raggedright\arraybackslash}p{0.2\linewidth}>{\raggedright\arraybackslash}p{0.4\linewidth}>{\raggedright\arraybackslash}p{0.36\linewidth}@{}}
\toprule
Model & Defaults & Validation grid \\
\midrule
Polynomial ridge & degree 3, $\alpha=10^{-4}$ & degree $\{2,3\}$, $\alpha\in\{10^{-6},10^{-4},10^{-2}\}$ \\
Random forest & 300 trees, all features per split, min.\ leaf 1 & max.\ features $\{0.33,0.66,1\}$, min.\ leaf $\{1,3\}$ \\
GBDT (LightGBM) & 1000 trees, lr 0.05, 31 leaves, subsample 0.8, colsample 0.8 & leaves $\{15,31,63\}$, lr $\{0.03,0.1\}$ \\
MLP & $128\times128$, $\alpha=10^{-4}$, Adam lr $10^{-3}$, batch 128, early stopping (10\,\% hold-out, patience 150) & layers $\{128^2, 128^3, 256^3\}$, $\alpha\in\{10^{-5},10^{-3}\}$ \\
\bottomrule
\end{tabular}
\caption{Model types, default hyperparameters and validation search grids. GBDT uses one
LightGBM regressor per current; the MLP also standardises its targets.}
\label{tab:hyper}
\end{table}

The MLP patience was raised from 30 to 150 epochs after we saw that early stopping ended
training after only 150--200 epochs; on validation the longer patience cut the
\texttt{nmos\_on} $I_D$/$I_S$ 95th-percentile error from 34/48\,\% to 24/24\,\% and the
\texttt{nmos\_off} $I_S$ error from 144\,\% to 87\,\%.

\subsection{Selection and metrics}
Errors are relative, $e = (I_{ML}-I_{SPICE})/|I_{SPICE}|$, and we report the median and 95th
percentile of $|e|$, plus $R^2$ on $\log_{10}|I|$. Samples with $|I_{SPICE}| < 10^{-15}$\,A are
excluded from the relative errors (they would dominate them while contributing nothing to
leakage); this removes all OFF-state gate currents ($\sim\!10^{-22}$\,A). For every case the
model with the lowest validation 95th-percentile error, averaged over the currents that have
one, is selected; ties are broken by the median error and then by training time.

\subsection{Results}
\todo{paste results/tables/model\_comparison\_summary.md (validation, 4 models $\times$ 4
cases: $R^2$, median, p95) as Table~X; the long per-current table goes to the appendix}

\todo{paste results/tables/chosen\_models.md: chosen model per case and its validation score}

\todo{paste results/tables/chosen\_models\_test.md: test median and p95 per case and current}

\todo{include results/plots/pred\_vs\_spice\_\{nmos\_on,nmos\_off,pmos\_on,pmos\_off\}.png as a
2$\times$2 figure}

\todo{one paragraph: which model type won per case and why (first round: GBDT for the ON
cases, MLP for the OFF cases; polynomial ridge was far worse everywhere)}


\section{Error analysis: transistor models}\label{sec:error-transistor}
% P3 (gates) and P4 (circuit) add their own subsections.

The following was measured on the \emph{validation} set with the first-round models (before
channel shaping); \todo{re-measure with the final models and update the numbers}.

\paragraph{ON-state $I_D$, $I_S$ near $V_D = V_S$.}
The largest relative errors by far are ON-state drain and source currents at small $|V_{DS}|$.
For \texttt{nmos\_on} $I_D$ the 95th-percentile error is 506\,\% for $|V_{DS}|<10$\,mV,
67\,\% for 30--100\,mV and 30\,\% above 0.3\,V (\texttt{pmos\_on}: 544\,\%, 66\,\%, 29\,\%).
The channel current passes through zero there, so a small absolute error is a large relative
one. This motivated the channel shaping above.

\paragraph{OFF-state source current.}
OFF-state $I_S$ is the worst-predicted current overall (95th-percentile error
$\approx 100$\,\% for $|V_{DS}|<0.1$\,V). It is tiny (often below $10^{-14}$\,A), changes
sign, and over 40\,\% of its samples fall below the $10^{-15}$\,A floor.

\paragraph{Temperature.}
OFF-state body (junction) current gets harder to predict as temperature rises: for
\texttt{nmos\_off} $I_B$ the 95th-percentile error grows from 6.5\,\% below 0\,$^\circ$C to
21\,\% above 85\,$^\circ$C, following the steep temperature dependence of junction leakage.

\paragraph{Coverage of the operating point.}
In a static CMOS gate every ON transistor conducts no DC current, so its drain and source sit
at the same rail: $V_{DS}\approx 0$ (in our gate simulations the median is about
$2\times10^{-7}$\,V). The first-round training data sampled $V_D$ and $V_S$ uniformly in
$[0,V_{DD}]$, so only 26 of 10\,000 \texttt{nmos\_on} rows and 17 of 10\,000 \texttt{pmos\_on}
rows had $|V_{DS}|<1$\,mV. The ON models were therefore asked to predict at a point they had
almost never seen, which is also where they are least accurate.
\todo{describe P1's low-$V_{DS}$ sampling and its effect once the data is regenerated}
(The gate data was used only to identify this problem; it was never used to train, tune or
select transistor models.)


\section{Limitations and improvements}\label{sec:limitations}
% P2 leads; P1, P3, P4 add their points.

\begin{itemize}
  \item \textbf{Training data vs.\ operating point.} Uniform sampling of $V_D$, $V_S$
        under-represents the bias points that matter in gates ($V_{DS}\approx 0$ for ON
        devices). Sampling should follow where the models will be used.
        \todo{update once low-$V_{DS}$ data is in}
  \item \textbf{Relative error near zero current.} Currents that cross zero (ON $I_D$, $I_S$)
        or are tiny (OFF $I_S$) have unbounded relative error; the $10^{-15}$\,A floor only
        partly addresses this. Their contribution to gate leakage is small, but they dominate
        the per-current statistics.
  \item \textbf{Gate voltage fixed at $0$ or $V_{DD}$.} The models only see ideal gate
        voltages; inside AND2/OR2 the internal inverter input is snapped to the nearest rail.
  \item \textbf{Independent models per current.} Kirchhoff's law ($\sum I = 0$) is not
        enforced; predicting three currents and deriving the fourth would guarantee it.
  \item \textbf{MLP training is sensitive to early stopping.} The patience setting alone
        changed the error by about a third; a learning-rate schedule or longer training could
        help further.
  \item \todo{learning-curve result with P1: does the error still fall at 7\,000 training
        rows? If yes, more data is an improvement.}
\end{itemize}
   from the main report (or pandoc it into report.md).
% Needs: \usepackage{amsmath, array, booktabs, graphicx, xcolor}
%        \newcommand{\todo}[1]{\textcolor{red}{[TODO: #1]}}
%
% Numbers marked "first-round" come from the models committed in 1f1d9f7 / validation
% experiments run on 2026-11-07. Everything that depends on the FINAL retrain (after the
% low-Vds data from P1) is marked \todo and must be pasted from results/tables/*.md.

\section{Transistor models}\label{sec:models}

\subsection{Data split}
Each of the four single-transistor datasets (\texttt{nmos\_on}, \texttt{nmos\_off},
\texttt{pmos\_on}, \texttt{pmos\_off}) is shuffled with a fixed seed (\texttt{seeds.split}~$=42$)
and split 70/15/15 into training, validation and test sets (7\,000/1\,500/1\,500 rows for
10\,000 samples). The row indices of every split are saved in \texttt{artifacts/splits/} so
the split is reproducible and auditable. All model choices---model type, hyperparameters and
preprocessing---are made on the validation set only; the test set is used exactly once, in
\texttt{scripts/05\_finalize\_models.py}, to report the chosen models.
\todo{update counts if P1's low-$V_{DS}$ sampling changes the dataset size}

\subsection{Input features}
The raw inputs are $V_{DD}$, $T$, $W$, $V_D$, $V_S$, $V_G$, $V_B$ and the five varied BSIM4
parameters (\texttt{vth0}, \texttt{toxe}, \texttt{ndep}, \texttt{u0}, \texttt{xj}). Within one
case $V_G$ and $V_B$ are always $0$ or $V_{DD}$, so they carry no information on their own; the
terminal voltage \emph{differences} do. From the raw inputs we build 22 features:
\begin{itemize}
  \item supply and temperature: $V_{DD}$, $T$ (K), $1/T$, thermal voltage $v_T = kT/q$;
  \item geometry: $W$, $\ln W$;
  \item bias: $V_{GS}$, $V_{DS}$, $V_{BS}$, $V_{GD}$, $V_{BD}$, $|V_{DS}|$;
  \item subthreshold physics: $V_{GS}/v_T$, $V_{DS}/v_T$, $(V_{GS}-|V_{th0}|)/v_T$ and the drain
        factor $1-e^{-|V_{DS}|/v_T}$;
  \item process: the five parameters divided by their nominal values from the model card
        (so they lie in $[0.9, 1.1]$), plus $1/t_{ox}$, since gate tunnelling grows steeply as
        the oxide gets thinner.
\end{itemize}
For PMOS every voltage difference is negated (e.g.\ $V_{SG}$ instead of $V_{GS}$) and
$|V_{th0}|$ is used, so PMOS features read like NMOS ones: $V_{GS}>0$ always means ``more on''.
The same function computes the features in training and inside the saved model, and the
feature list is stored in the model file; a model refuses to predict if the feature code
has changed since it was trained.

\subsection{Target transform}
The four terminal currents span about ten decades (from $\sim\!10^{-22}$\,A for OFF-state
gate current to $\sim\!10^{-4}$\,A for ON-state drain current) and $I_D$, $I_S$ change sign
with $V_D-V_S$. Regressing amperes with a squared loss would fit only the largest currents, so
each model predicts
\begin{equation}
  z_t = \operatorname{asinh}\!\left(\frac{I_t}{i_{0,t}}\right),\qquad
  t \in \{I_D, I_G, I_S, I_B\},
\end{equation}
where $i_{0,t}$ is the 1st percentile of $|I_t|$ on the training rows. For $|I|\gg i_0$,
$\operatorname{asinh}(I/i_0)\approx \operatorname{sign}(I)\ln(2|I|/i_0)$, so an error
$\delta z$ is a relative current error of about $\delta z$, which is exactly what we are judged
on. Unlike $\log_{10}|I|$, asinh is smooth through zero and carries the sign, so no separate sign
model is needed. A per-target $i_0$ is necessary because the currents live on very different
scales: one fixed $i_0$ cannot be ``small'' for all of them.

\paragraph{Channel shaping (ON cases).}
In the ON cases the channel current passes through zero at $V_D=V_S$ with the shape
$1-e^{-V_{DS}/v_T}$ in subthreshold and $\propto V_{DS}$ in the linear region. For
\texttt{nmos\_on} and \texttt{pmos\_on}, $I_D$ and $I_S$ are therefore divided by
\begin{equation}
  s(V_{DS}) = \tanh\!\left(\frac{V_D - V_S}{2 v_T}\right)
\end{equation}
before the transform and multiplied back afterwards; the model then learns a smooth,
one-signed quantity. On the validation set this reduced the 95th-percentile error of the
\texttt{nmos\_on} MLP from 164\,\% / 206\,\% to 34\,\% / 48\,\% for $I_D$ / $I_S$. It is
\emph{not} used for the OFF cases: there $I_D$ and $I_S$ near $V_D=V_S$ are dominated by
junction leakage, which does not change sign with $V_{DS}$, and shaping made them worse
(\texttt{nmos\_off} GBDT: 65\,\% / 100\,\% $\rightarrow$ 70\,\% / 111\,\%).

\subsection{Model types and hyperparameters}
Four model types are compared for every case (Table~\ref{tab:hyper}). Each predicts all four
transformed currents at once. Inputs are standardised, and every estimator is seeded
(\texttt{seeds.training}~$=7$). Each model type has a small grid (6 points) that is searched
on the validation set; the grid point with the lowest validation score is kept.

\begin{table}[h]
\centering\small
\begin{tabular}{@{}>{\raggedright\arraybackslash}p{0.2\linewidth}>{\raggedright\arraybackslash}p{0.4\linewidth}>{\raggedright\arraybackslash}p{0.36\linewidth}@{}}
\toprule
Model & Defaults & Validation grid \\
\midrule
Polynomial ridge & degree 3, $\alpha=10^{-4}$ & degree $\{2,3\}$, $\alpha\in\{10^{-6},10^{-4},10^{-2}\}$ \\
Random forest & 300 trees, all features per split, min.\ leaf 1 & max.\ features $\{0.33,0.66,1\}$, min.\ leaf $\{1,3\}$ \\
GBDT (LightGBM) & 1000 trees, lr 0.05, 31 leaves, subsample 0.8, colsample 0.8 & leaves $\{15,31,63\}$, lr $\{0.03,0.1\}$ \\
MLP & $128\times128$, $\alpha=10^{-4}$, Adam lr $10^{-3}$, batch 128, early stopping (10\,\% hold-out, patience 150) & layers $\{128^2, 128^3, 256^3\}$, $\alpha\in\{10^{-5},10^{-3}\}$ \\
\bottomrule
\end{tabular}
\caption{Model types, default hyperparameters and validation search grids. GBDT uses one
LightGBM regressor per current; the MLP also standardises its targets.}
\label{tab:hyper}
\end{table}

The MLP patience was raised from 30 to 150 epochs after we saw that early stopping ended
training after only 150--200 epochs; on validation the longer patience cut the
\texttt{nmos\_on} $I_D$/$I_S$ 95th-percentile error from 34/48\,\% to 24/24\,\% and the
\texttt{nmos\_off} $I_S$ error from 144\,\% to 87\,\%.

\subsection{Selection and metrics}
Errors are relative, $e = (I_{ML}-I_{SPICE})/|I_{SPICE}|$, and we report the median and 95th
percentile of $|e|$, plus $R^2$ on $\log_{10}|I|$. Samples with $|I_{SPICE}| < 10^{-15}$\,A are
excluded from the relative errors (they would dominate them while contributing nothing to
leakage); this removes all OFF-state gate currents ($\sim\!10^{-22}$\,A). For every case the
model with the lowest validation 95th-percentile error, averaged over the currents that have
one, is selected; ties are broken by the median error and then by training time.

\subsection{Results}
\todo{paste results/tables/model\_comparison\_summary.md (validation, 4 models $\times$ 4
cases: $R^2$, median, p95) as Table~X; the long per-current table goes to the appendix}

\todo{paste results/tables/chosen\_models.md: chosen model per case and its validation score}

\todo{paste results/tables/chosen\_models\_test.md: test median and p95 per case and current}

\todo{include results/plots/pred\_vs\_spice\_\{nmos\_on,nmos\_off,pmos\_on,pmos\_off\}.png as a
2$\times$2 figure}

\todo{one paragraph: which model type won per case and why (first round: GBDT for the ON
cases, MLP for the OFF cases; polynomial ridge was far worse everywhere)}


\section{Error analysis: transistor models}\label{sec:error-transistor}
% P3 (gates) and P4 (circuit) add their own subsections.

The following was measured on the \emph{validation} set with the first-round models (before
channel shaping); \todo{re-measure with the final models and update the numbers}.

\paragraph{ON-state $I_D$, $I_S$ near $V_D = V_S$.}
The largest relative errors by far are ON-state drain and source currents at small $|V_{DS}|$.
For \texttt{nmos\_on} $I_D$ the 95th-percentile error is 506\,\% for $|V_{DS}|<10$\,mV,
67\,\% for 30--100\,mV and 30\,\% above 0.3\,V (\texttt{pmos\_on}: 544\,\%, 66\,\%, 29\,\%).
The channel current passes through zero there, so a small absolute error is a large relative
one. This motivated the channel shaping above.

\paragraph{OFF-state source current.}
OFF-state $I_S$ is the worst-predicted current overall (95th-percentile error
$\approx 100$\,\% for $|V_{DS}|<0.1$\,V). It is tiny (often below $10^{-14}$\,A), changes
sign, and over 40\,\% of its samples fall below the $10^{-15}$\,A floor.

\paragraph{Temperature.}
OFF-state body (junction) current gets harder to predict as temperature rises: for
\texttt{nmos\_off} $I_B$ the 95th-percentile error grows from 6.5\,\% below 0\,$^\circ$C to
21\,\% above 85\,$^\circ$C, following the steep temperature dependence of junction leakage.

\paragraph{Coverage of the operating point.}
In a static CMOS gate every ON transistor conducts no DC current, so its drain and source sit
at the same rail: $V_{DS}\approx 0$ (in our gate simulations the median is about
$2\times10^{-7}$\,V). The first-round training data sampled $V_D$ and $V_S$ uniformly in
$[0,V_{DD}]$, so only 26 of 10\,000 \texttt{nmos\_on} rows and 17 of 10\,000 \texttt{pmos\_on}
rows had $|V_{DS}|<1$\,mV. The ON models were therefore asked to predict at a point they had
almost never seen, which is also where they are least accurate.
\todo{describe P1's low-$V_{DS}$ sampling and its effect once the data is regenerated}
(The gate data was used only to identify this problem; it was never used to train, tune or
select transistor models.)


\section{Limitations and improvements}\label{sec:limitations}
% P2 leads; P1, P3, P4 add their points.

\begin{itemize}
  \item \textbf{Training data vs.\ operating point.} Uniform sampling of $V_D$, $V_S$
        under-represents the bias points that matter in gates ($V_{DS}\approx 0$ for ON
        devices). Sampling should follow where the models will be used.
        \todo{update once low-$V_{DS}$ data is in}
  \item \textbf{Relative error near zero current.} Currents that cross zero (ON $I_D$, $I_S$)
        or are tiny (OFF $I_S$) have unbounded relative error; the $10^{-15}$\,A floor only
        partly addresses this. Their contribution to gate leakage is small, but they dominate
        the per-current statistics.
  \item \textbf{Gate voltage fixed at $0$ or $V_{DD}$.} The models only see ideal gate
        voltages; inside AND2/OR2 the internal inverter input is snapped to the nearest rail.
  \item \textbf{Independent models per current.} Kirchhoff's law ($\sum I = 0$) is not
        enforced; predicting three currents and deriving the fourth would guarantee it.
  \item \textbf{MLP training is sensitive to early stopping.} The patience setting alone
        changed the error by about a third; a learning-rate schedule or longer training could
        help further.
  \item \todo{learning-curve result with P1: does the error still fall at 7\,000 training
        rows? If yes, more data is an improvement.}
\end{itemize}
   from the main report (or pandoc it into report.md).
% Needs: \usepackage{amsmath, array, booktabs, graphicx, xcolor}
%        \newcommand{\todo}[1]{\textcolor{red}{[TODO: #1]}}
%
% Numbers marked "first-round" come from the models committed in 1f1d9f7 / validation
% experiments run on 2026-11-07. Everything that depends on the FINAL retrain (after the
% low-Vds data from P1) is marked \todo and must be pasted from results/tables/*.md.

\section{Transistor models}\label{sec:models}

\subsection{Data split}
Each of the four single-transistor datasets (\texttt{nmos\_on}, \texttt{nmos\_off},
\texttt{pmos\_on}, \texttt{pmos\_off}) is shuffled with a fixed seed (\texttt{seeds.split}~$=42$)
and split 70/15/15 into training, validation and test sets (7\,000/1\,500/1\,500 rows for
10\,000 samples). The row indices of every split are saved in \texttt{artifacts/splits/} so
the split is reproducible and auditable. All model choices---model type, hyperparameters and
preprocessing---are made on the validation set only; the test set is used exactly once, in
\texttt{scripts/05\_finalize\_models.py}, to report the chosen models.
\todo{update counts if P1's low-$V_{DS}$ sampling changes the dataset size}

\subsection{Input features}
The raw inputs are $V_{DD}$, $T$, $W$, $V_D$, $V_S$, $V_G$, $V_B$ and the five varied BSIM4
parameters (\texttt{vth0}, \texttt{toxe}, \texttt{ndep}, \texttt{u0}, \texttt{xj}). Within one
case $V_G$ and $V_B$ are always $0$ or $V_{DD}$, so they carry no information on their own; the
terminal voltage \emph{differences} do. From the raw inputs we build 22 features:
\begin{itemize}
  \item supply and temperature: $V_{DD}$, $T$ (K), $1/T$, thermal voltage $v_T = kT/q$;
  \item geometry: $W$, $\ln W$;
  \item bias: $V_{GS}$, $V_{DS}$, $V_{BS}$, $V_{GD}$, $V_{BD}$, $|V_{DS}|$;
  \item subthreshold physics: $V_{GS}/v_T$, $V_{DS}/v_T$, $(V_{GS}-|V_{th0}|)/v_T$ and the drain
        factor $1-e^{-|V_{DS}|/v_T}$;
  \item process: the five parameters divided by their nominal values from the model card
        (so they lie in $[0.9, 1.1]$), plus $1/t_{ox}$, since gate tunnelling grows steeply as
        the oxide gets thinner.
\end{itemize}
For PMOS every voltage difference is negated (e.g.\ $V_{SG}$ instead of $V_{GS}$) and
$|V_{th0}|$ is used, so PMOS features read like NMOS ones: $V_{GS}>0$ always means ``more on''.
The same function computes the features in training and inside the saved model, and the
feature list is stored in the model file; a model refuses to predict if the feature code
has changed since it was trained.

\subsection{Target transform}
The four terminal currents span about ten decades (from $\sim\!10^{-22}$\,A for OFF-state
gate current to $\sim\!10^{-4}$\,A for ON-state drain current) and $I_D$, $I_S$ change sign
with $V_D-V_S$. Regressing amperes with a squared loss would fit only the largest currents, so
each model predicts
\begin{equation}
  z_t = \operatorname{asinh}\!\left(\frac{I_t}{i_{0,t}}\right),\qquad
  t \in \{I_D, I_G, I_S, I_B\},
\end{equation}
where $i_{0,t}$ is the 1st percentile of $|I_t|$ on the training rows. For $|I|\gg i_0$,
$\operatorname{asinh}(I/i_0)\approx \operatorname{sign}(I)\ln(2|I|/i_0)$, so an error
$\delta z$ is a relative current error of about $\delta z$, which is exactly what we are judged
on. Unlike $\log_{10}|I|$, asinh is smooth through zero and carries the sign, so no separate sign
model is needed. A per-target $i_0$ is necessary because the currents live on very different
scales: one fixed $i_0$ cannot be ``small'' for all of them.

\paragraph{Channel shaping (ON cases).}
In the ON cases the channel current passes through zero at $V_D=V_S$ with the shape
$1-e^{-V_{DS}/v_T}$ in subthreshold and $\propto V_{DS}$ in the linear region. For
\texttt{nmos\_on} and \texttt{pmos\_on}, $I_D$ and $I_S$ are therefore divided by
\begin{equation}
  s(V_{DS}) = \tanh\!\left(\frac{V_D - V_S}{2 v_T}\right)
\end{equation}
before the transform and multiplied back afterwards; the model then learns a smooth,
one-signed quantity. On the validation set this reduced the 95th-percentile error of the
\texttt{nmos\_on} MLP from 164\,\% / 206\,\% to 34\,\% / 48\,\% for $I_D$ / $I_S$. It is
\emph{not} used for the OFF cases: there $I_D$ and $I_S$ near $V_D=V_S$ are dominated by
junction leakage, which does not change sign with $V_{DS}$, and shaping made them worse
(\texttt{nmos\_off} GBDT: 65\,\% / 100\,\% $\rightarrow$ 70\,\% / 111\,\%).

\subsection{Model types and hyperparameters}
Four model types are compared for every case (Table~\ref{tab:hyper}). Each predicts all four
transformed currents at once. Inputs are standardised, and every estimator is seeded
(\texttt{seeds.training}~$=7$). Each model type has a small grid (6 points) that is searched
on the validation set; the grid point with the lowest validation score is kept.

\begin{table}[h]
\centering\small
\begin{tabular}{@{}>{\raggedright\arraybackslash}p{0.2\linewidth}>{\raggedright\arraybackslash}p{0.4\linewidth}>{\raggedright\arraybackslash}p{0.36\linewidth}@{}}
\toprule
Model & Defaults & Validation grid \\
\midrule
Polynomial ridge & degree 3, $\alpha=10^{-4}$ & degree $\{2,3\}$, $\alpha\in\{10^{-6},10^{-4},10^{-2}\}$ \\
Random forest & 300 trees, all features per split, min.\ leaf 1 & max.\ features $\{0.33,0.66,1\}$, min.\ leaf $\{1,3\}$ \\
GBDT (LightGBM) & 1000 trees, lr 0.05, 31 leaves, subsample 0.8, colsample 0.8 & leaves $\{15,31,63\}$, lr $\{0.03,0.1\}$ \\
MLP & $128\times128$, $\alpha=10^{-4}$, Adam lr $10^{-3}$, batch 128, early stopping (10\,\% hold-out, patience 150) & layers $\{128^2, 128^3, 256^3\}$, $\alpha\in\{10^{-5},10^{-3}\}$ \\
\bottomrule
\end{tabular}
\caption{Model types, default hyperparameters and validation search grids. GBDT uses one
LightGBM regressor per current; the MLP also standardises its targets.}
\label{tab:hyper}
\end{table}

The MLP patience was raised from 30 to 150 epochs after we saw that early stopping ended
training after only 150--200 epochs; on validation the longer patience cut the
\texttt{nmos\_on} $I_D$/$I_S$ 95th-percentile error from 34/48\,\% to 24/24\,\% and the
\texttt{nmos\_off} $I_S$ error from 144\,\% to 87\,\%.

\subsection{Selection and metrics}
Errors are relative, $e = (I_{ML}-I_{SPICE})/|I_{SPICE}|$, and we report the median and 95th
percentile of $|e|$, plus $R^2$ on $\log_{10}|I|$. Samples with $|I_{SPICE}| < 10^{-15}$\,A are
excluded from the relative errors (they would dominate them while contributing nothing to
leakage); this removes all OFF-state gate currents ($\sim\!10^{-22}$\,A). For every case the
model with the lowest validation 95th-percentile error, averaged over the currents that have
one, is selected; ties are broken by the median error and then by training time.

\subsection{Results}
\todo{paste results/tables/model\_comparison\_summary.md (validation, 4 models $\times$ 4
cases: $R^2$, median, p95) as Table~X; the long per-current table goes to the appendix}

\todo{paste results/tables/chosen\_models.md: chosen model per case and its validation score}

\todo{paste results/tables/chosen\_models\_test.md: test median and p95 per case and current}

\todo{include results/plots/pred\_vs\_spice\_\{nmos\_on,nmos\_off,pmos\_on,pmos\_off\}.png as a
2$\times$2 figure}

\todo{one paragraph: which model type won per case and why (first round: GBDT for the ON
cases, MLP for the OFF cases; polynomial ridge was far worse everywhere)}


\section{Error analysis: transistor models}\label{sec:error-transistor}
% P3 (gates) and P4 (circuit) add their own subsections.

The following was measured on the \emph{validation} set with the first-round models (before
channel shaping); \todo{re-measure with the final models and update the numbers}.

\paragraph{ON-state $I_D$, $I_S$ near $V_D = V_S$.}
The largest relative errors by far are ON-state drain and source currents at small $|V_{DS}|$.
For \texttt{nmos\_on} $I_D$ the 95th-percentile error is 506\,\% for $|V_{DS}|<10$\,mV,
67\,\% for 30--100\,mV and 30\,\% above 0.3\,V (\texttt{pmos\_on}: 544\,\%, 66\,\%, 29\,\%).
The channel current passes through zero there, so a small absolute error is a large relative
one. This motivated the channel shaping above.

\paragraph{OFF-state source current.}
OFF-state $I_S$ is the worst-predicted current overall (95th-percentile error
$\approx 100$\,\% for $|V_{DS}|<0.1$\,V). It is tiny (often below $10^{-14}$\,A), changes
sign, and over 40\,\% of its samples fall below the $10^{-15}$\,A floor.

\paragraph{Temperature.}
OFF-state body (junction) current gets harder to predict as temperature rises: for
\texttt{nmos\_off} $I_B$ the 95th-percentile error grows from 6.5\,\% below 0\,$^\circ$C to
21\,\% above 85\,$^\circ$C, following the steep temperature dependence of junction leakage.

\paragraph{Coverage of the operating point.}
In a static CMOS gate every ON transistor conducts no DC current, so its drain and source sit
at the same rail: $V_{DS}\approx 0$ (in our gate simulations the median is about
$2\times10^{-7}$\,V). The first-round training data sampled $V_D$ and $V_S$ uniformly in
$[0,V_{DD}]$, so only 26 of 10\,000 \texttt{nmos\_on} rows and 17 of 10\,000 \texttt{pmos\_on}
rows had $|V_{DS}|<1$\,mV. The ON models were therefore asked to predict at a point they had
almost never seen, which is also where they are least accurate.
\todo{describe P1's low-$V_{DS}$ sampling and its effect once the data is regenerated}
(The gate data was used only to identify this problem; it was never used to train, tune or
select transistor models.)


\section{Limitations and improvements}\label{sec:limitations}
% P2 leads; P1, P3, P4 add their points.

\begin{itemize}
  \item \textbf{Training data vs.\ operating point.} Uniform sampling of $V_D$, $V_S$
        under-represents the bias points that matter in gates ($V_{DS}\approx 0$ for ON
        devices). Sampling should follow where the models will be used.
        \todo{update once low-$V_{DS}$ data is in}
  \item \textbf{Relative error near zero current.} Currents that cross zero (ON $I_D$, $I_S$)
        or are tiny (OFF $I_S$) have unbounded relative error; the $10^{-15}$\,A floor only
        partly addresses this. Their contribution to gate leakage is small, but they dominate
        the per-current statistics.
  \item \textbf{Gate voltage fixed at $0$ or $V_{DD}$.} The models only see ideal gate
        voltages; inside AND2/OR2 the internal inverter input is snapped to the nearest rail.
  \item \textbf{Independent models per current.} Kirchhoff's law ($\sum I = 0$) is not
        enforced; predicting three currents and deriving the fourth would guarantee it.
  \item \textbf{MLP training is sensitive to early stopping.} The patience setting alone
        changed the error by about a third; a learning-rate schedule or longer training could
        help further.
  \item \todo{learning-curve result with P1: does the error still fall at 7\,000 training
        rows? If yes, more data is an improvement.}
\end{itemize}
   from the main report (or pandoc it into report.md).
% Needs: \usepackage{amsmath, array, booktabs, graphicx, xcolor}
%        \newcommand{\todo}[1]{\textcolor{red}{[TODO: #1]}}
%
% Numbers marked "first-round" come from the models committed in 1f1d9f7 / validation
% experiments run on 2026-11-07. Everything that depends on the FINAL retrain (after the
% low-Vds data from P1) is marked \todo and must be pasted from results/tables/*.md.

\section{Transistor models}\label{sec:models}

\subsection{Data split}
Each of the four single-transistor datasets (\texttt{nmos\_on}, \texttt{nmos\_off},
\texttt{pmos\_on}, \texttt{pmos\_off}) is shuffled with a fixed seed (\texttt{seeds.split}~$=42$)
and split 70/15/15 into training, validation and test sets (7\,000/1\,500/1\,500 rows for
10\,000 samples). The row indices of every split are saved in \texttt{artifacts/splits/} so
the split is reproducible and auditable. All model choices---model type, hyperparameters and
preprocessing---are made on the validation set only; the test set is used exactly once, in
\texttt{scripts/05\_finalize\_models.py}, to report the chosen models.
\todo{update counts if P1's low-$V_{DS}$ sampling changes the dataset size}

\subsection{Input features}
The raw inputs are $V_{DD}$, $T$, $W$, $V_D$, $V_S$, $V_G$, $V_B$ and the five varied BSIM4
parameters (\texttt{vth0}, \texttt{toxe}, \texttt{ndep}, \texttt{u0}, \texttt{xj}). Within one
case $V_G$ and $V_B$ are always $0$ or $V_{DD}$, so they carry no information on their own; the
terminal voltage \emph{differences} do. From the raw inputs we build 22 features:
\begin{itemize}
  \item supply and temperature: $V_{DD}$, $T$ (K), $1/T$, thermal voltage $v_T = kT/q$;
  \item geometry: $W$, $\ln W$;
  \item bias: $V_{GS}$, $V_{DS}$, $V_{BS}$, $V_{GD}$, $V_{BD}$, $|V_{DS}|$;
  \item subthreshold physics: $V_{GS}/v_T$, $V_{DS}/v_T$, $(V_{GS}-|V_{th0}|)/v_T$ and the drain
        factor $1-e^{-|V_{DS}|/v_T}$;
  \item process: the five parameters divided by their nominal values from the model card
        (so they lie in $[0.9, 1.1]$), plus $1/t_{ox}$, since gate tunnelling grows steeply as
        the oxide gets thinner.
\end{itemize}
For PMOS every voltage difference is negated (e.g.\ $V_{SG}$ instead of $V_{GS}$) and
$|V_{th0}|$ is used, so PMOS features read like NMOS ones: $V_{GS}>0$ always means ``more on''.
The same function computes the features in training and inside the saved model, and the
feature list is stored in the model file; a model refuses to predict if the feature code
has changed since it was trained.

\subsection{Target transform}
The four terminal currents span about ten decades (from $\sim\!10^{-22}$\,A for OFF-state
gate current to $\sim\!10^{-4}$\,A for ON-state drain current) and $I_D$, $I_S$ change sign
with $V_D-V_S$. Regressing amperes with a squared loss would fit only the largest currents, so
each model predicts
\begin{equation}
  z_t = \operatorname{asinh}\!\left(\frac{I_t}{i_{0,t}}\right),\qquad
  t \in \{I_D, I_G, I_S, I_B\},
\end{equation}
where $i_{0,t}$ is the 1st percentile of $|I_t|$ on the training rows. For $|I|\gg i_0$,
$\operatorname{asinh}(I/i_0)\approx \operatorname{sign}(I)\ln(2|I|/i_0)$, so an error
$\delta z$ is a relative current error of about $\delta z$, which is exactly what we are judged
on. Unlike $\log_{10}|I|$, asinh is smooth through zero and carries the sign, so no separate sign
model is needed. A per-target $i_0$ is necessary because the currents live on very different
scales: one fixed $i_0$ cannot be ``small'' for all of them.

\paragraph{Channel shaping (ON cases).}
In the ON cases the channel current passes through zero at $V_D=V_S$ with the shape
$1-e^{-V_{DS}/v_T}$ in subthreshold and $\propto V_{DS}$ in the linear region. For
\texttt{nmos\_on} and \texttt{pmos\_on}, $I_D$ and $I_S$ are therefore divided by
\begin{equation}
  s(V_{DS}) = \tanh\!\left(\frac{V_D - V_S}{2 v_T}\right)
\end{equation}
before the transform and multiplied back afterwards; the model then learns a smooth,
one-signed quantity. On the validation set this reduced the 95th-percentile error of the
\texttt{nmos\_on} MLP from 164\,\% / 206\,\% to 34\,\% / 48\,\% for $I_D$ / $I_S$. It is
\emph{not} used for the OFF cases: there $I_D$ and $I_S$ near $V_D=V_S$ are dominated by
junction leakage, which does not change sign with $V_{DS}$, and shaping made them worse
(\texttt{nmos\_off} GBDT: 65\,\% / 100\,\% $\rightarrow$ 70\,\% / 111\,\%).

\subsection{Model types and hyperparameters}
Four model types are compared for every case (Table~\ref{tab:hyper}). Each predicts all four
transformed currents at once. Inputs are standardised, and every estimator is seeded
(\texttt{seeds.training}~$=7$). Each model type has a small grid (6 points) that is searched
on the validation set; the grid point with the lowest validation score is kept.

\begin{table}[h]
\centering\small
\begin{tabular}{@{}>{\raggedright\arraybackslash}p{0.2\linewidth}>{\raggedright\arraybackslash}p{0.4\linewidth}>{\raggedright\arraybackslash}p{0.36\linewidth}@{}}
\toprule
Model & Defaults & Validation grid \\
\midrule
Polynomial ridge & degree 3, $\alpha=10^{-4}$ & degree $\{2,3\}$, $\alpha\in\{10^{-6},10^{-4},10^{-2}\}$ \\
Random forest & 300 trees, all features per split, min.\ leaf 1 & max.\ features $\{0.33,0.66,1\}$, min.\ leaf $\{1,3\}$ \\
GBDT (LightGBM) & 1000 trees, lr 0.05, 31 leaves, subsample 0.8, colsample 0.8 & leaves $\{15,31,63\}$, lr $\{0.03,0.1\}$ \\
MLP & $128\times128$, $\alpha=10^{-4}$, Adam lr $10^{-3}$, batch 128, early stopping (10\,\% hold-out, patience 150) & layers $\{128^2, 128^3, 256^3\}$, $\alpha\in\{10^{-5},10^{-3}\}$ \\
\bottomrule
\end{tabular}
\caption{Model types, default hyperparameters and validation search grids. GBDT uses one
LightGBM regressor per current; the MLP also standardises its targets.}
\label{tab:hyper}
\end{table}

The MLP patience was raised from 30 to 150 epochs after we saw that early stopping ended
training after only 150--200 epochs; on validation the longer patience cut the
\texttt{nmos\_on} $I_D$/$I_S$ 95th-percentile error from 34/48\,\% to 24/24\,\% and the
\texttt{nmos\_off} $I_S$ error from 144\,\% to 87\,\%.

\subsection{Selection and metrics}
Errors are relative, $e = (I_{ML}-I_{SPICE})/|I_{SPICE}|$, and we report the median and 95th
percentile of $|e|$, plus $R^2$ on $\log_{10}|I|$. Samples with $|I_{SPICE}| < 10^{-15}$\,A are
excluded from the relative errors (they would dominate them while contributing nothing to
leakage); this removes all OFF-state gate currents ($\sim\!10^{-22}$\,A). For every case the
model with the lowest validation 95th-percentile error, averaged over the currents that have
one, is selected; ties are broken by the median error and then by training time.

\subsection{Results}
\todo{paste results/tables/model\_comparison\_summary.md (validation, 4 models $\times$ 4
cases: $R^2$, median, p95) as Table~X; the long per-current table goes to the appendix}

\todo{paste results/tables/chosen\_models.md: chosen model per case and its validation score}

\todo{paste results/tables/chosen\_models\_test.md: test median and p95 per case and current}

\todo{include results/plots/pred\_vs\_spice\_\{nmos\_on,nmos\_off,pmos\_on,pmos\_off\}.png as a
2$\times$2 figure}

\todo{one paragraph: which model type won per case and why (first round: GBDT for the ON
cases, MLP for the OFF cases; polynomial ridge was far worse everywhere)}


\section{Error analysis: transistor models}\label{sec:error-transistor}
% P3 (gates) and P4 (circuit) add their own subsections.

The following was measured on the \emph{validation} set with the first-round models (before
channel shaping); \todo{re-measure with the final models and update the numbers}.

\paragraph{ON-state $I_D$, $I_S$ near $V_D = V_S$.}
The largest relative errors by far are ON-state drain and source currents at small $|V_{DS}|$.
For \texttt{nmos\_on} $I_D$ the 95th-percentile error is 506\,\% for $|V_{DS}|<10$\,mV,
67\,\% for 30--100\,mV and 30\,\% above 0.3\,V (\texttt{pmos\_on}: 544\,\%, 66\,\%, 29\,\%).
The channel current passes through zero there, so a small absolute error is a large relative
one. This motivated the channel shaping above.

\paragraph{OFF-state source current.}
OFF-state $I_S$ is the worst-predicted current overall (95th-percentile error
$\approx 100$\,\% for $|V_{DS}|<0.1$\,V). It is tiny (often below $10^{-14}$\,A), changes
sign, and over 40\,\% of its samples fall below the $10^{-15}$\,A floor.

\paragraph{Temperature.}
OFF-state body (junction) current gets harder to predict as temperature rises: for
\texttt{nmos\_off} $I_B$ the 95th-percentile error grows from 6.5\,\% below 0\,$^\circ$C to
21\,\% above 85\,$^\circ$C, following the steep temperature dependence of junction leakage.

\paragraph{Coverage of the operating point.}
In a static CMOS gate every ON transistor conducts no DC current, so its drain and source sit
at the same rail: $V_{DS}\approx 0$ (in our gate simulations the median is about
$2\times10^{-7}$\,V). The first-round training data sampled $V_D$ and $V_S$ uniformly in
$[0,V_{DD}]$, so only 26 of 10\,000 \texttt{nmos\_on} rows and 17 of 10\,000 \texttt{pmos\_on}
rows had $|V_{DS}|<1$\,mV. The ON models were therefore asked to predict at a point they had
almost never seen, which is also where they are least accurate.
\todo{describe P1's low-$V_{DS}$ sampling and its effect once the data is regenerated}
(The gate data was used only to identify this problem; it was never used to train, tune or
select transistor models.)


\section{Limitations and improvements}\label{sec:limitations}
% P2 leads; P1, P3, P4 add their points.

\begin{itemize}
  \item \textbf{Training data vs.\ operating point.} Uniform sampling of $V_D$, $V_S$
        under-represents the bias points that matter in gates ($V_{DS}\approx 0$ for ON
        devices). Sampling should follow where the models will be used.
        \todo{update once low-$V_{DS}$ data is in}
  \item \textbf{Relative error near zero current.} Currents that cross zero (ON $I_D$, $I_S$)
        or are tiny (OFF $I_S$) have unbounded relative error; the $10^{-15}$\,A floor only
        partly addresses this. Their contribution to gate leakage is small, but they dominate
        the per-current statistics.
  \item \textbf{Gate voltage fixed at $0$ or $V_{DD}$.} The models only see ideal gate
        voltages; inside AND2/OR2 the internal inverter input is snapped to the nearest rail.
  \item \textbf{Independent models per current.} Kirchhoff's law ($\sum I = 0$) is not
        enforced; predicting three currents and deriving the fourth would guarantee it.
  \item \textbf{MLP training is sensitive to early stopping.} The patience setting alone
        changed the error by about a third; a learning-rate schedule or longer training could
        help further.
  \item \todo{learning-curve result with P1: does the error still fall at 7\,000 training
        rows? If yes, more data is an improvement.}
\end{itemize}
